% !TEX root = 第三章_矩阵(答案版).tex
\setcounter{chapter}{2}
\pagestyle{fancy}

\chapter{矩阵}

\section{进阶知识补充}

\subsection{迹}

\begin{definition}
方阵 $A=(a_{ij})\in M_n(F)$ 的迹(trace)定义为主对角线元素之和:
\[
\tr(A)=a_{11}+a_{22}+\cdots+a_{nn}.
\]
\end{definition}

\begin{theorem}[\markstar]
设 $A,B\in M_n(F)$,$\lambda\in F$,则:
\begin{enumerate}
  \item $\tr(A+B)=\tr(A)+\tr(B)$;
  \item $\tr(\lambda A)=\lambda\tr(A)$;
  \item $\tr(A)=\tr(A^{\mathsf T})$;
  \item $\tr(AB)=\tr(BA)$;
  \item 若 $F\subseteq\mathbb{R}$,则
  \[
  \tr(AA^{\mathsf T})=0\iff A=O.
  \]
  复数情形应将转置改为共轭转置,即
  \[
  \tr(AA^{\mathrm H})=0\iff A=O.
  \]
\end{enumerate}
\end{theorem}

\begin{remark}
复数情形若不加共轭,结论一般不成立.例如
\[
A=\begin{pmatrix}1&i\\ i&1\end{pmatrix}
\]
满足 $\tr(AA^{\mathsf T})=0$,但 $A\neq O$.性质 (5) 常用于证明矩阵等式:若要证明 $A=B$,可以尝试证明
\[
\tr\left((A-B)(A-B)^{\mathsf T}\right)=0.
\]
\end{remark}

%例题3.1.1
\begin{example}
设 $A$ 为二阶方阵,且存在方阵 $B$ 使得
\[
AB-BA=A.
\]
证明:$A^2=O$.
\end{example}
\exampleblank

\begin{solution}
由
\[
AB-BA=A
\]
取迹得 $\tr(A)=0$.进一步,利用
\[
A^2B-BA^2=A(AB-BA)+(AB-BA)A=2A^2,
\]
取迹得 $\tr(A^2)=0$.

设 $A$ 的两个特征值为 $\lambda_1,\lambda_2$,则
\[
\lambda_1+\lambda_2=0,
\qquad
\lambda_1^2+\lambda_2^2=0,
\]
从而 $\lambda_1=\lambda_2=0$.故 $A$ 的特征多项式为 $x^2$.由 Cayley--Hamilton 定理,
\[
\boxed{A^2=O}.
\]
\end{solution}

%例题3.1.2
\begin{example}
证明:对任意二阶复方阵 $A,B,C$,有
\[
A(BC-CB)^2-(BC-CB)^2A=O.
\]
\end{example}
\exampleblank

\begin{solution}
令
\[
D=BC-CB.
\]
由迹的循环性质,$\tr(D)=0$.对二阶矩阵 $D$ 使用 Cayley--Hamilton 定理:
\[
D^2-\tr(D)D+|D|I_2=O,
\]
故
\[
D^2=-|D|I_2.
\]
因此 $D^2$ 为数量矩阵,与任意二阶矩阵 $A$ 可交换,故
\[
AD^2-D^2A=O.
\]
\end{solution}

%例题3.1.3
\begin{example}
证明:对任意 $n$ 阶方阵 $A,B$,等式
\[
AB-BA=I_n
\]
都不成立.
\end{example}
\exampleblank

\begin{solution}
若等式成立,则两边取迹:
\[
\tr(AB)-\tr(BA)=\tr(I_n)=n.
\]
左端由 $\tr(AB)=\tr(BA)$ 等于 $0$,于是得到 $0=n$,矛盾.因此该等式不可能成立.
\end{solution}

%例题3.1.4
\begin{example}
设 $A$ 是 $n$ 阶实方阵,且
\[
A^2=AA^{\mathsf T}.
\]
证明:$A$ 是对称矩阵.
\end{example}
\exampleblank

\begin{solution}
对已知等式取转置,得
\[
(A^{\mathsf T})^2=AA^{\mathsf T}.
\]
令
\[
D=A-A^{\mathsf T}.
\]
则
\[
DD^{\mathsf T}
=(A-A^{\mathsf T})(A^{\mathsf T}-A)
=AA^{\mathsf T}-A^2-(A^{\mathsf T})^2+A^{\mathsf T}A.
\]
由两个已知关系,
\[
\tr(DD^{\mathsf T})
=\tr(A^{\mathsf T}A)-\tr(AA^{\mathsf T})=0.
\]
因 $A$ 为实矩阵,由迹的性质知 $D=O$,即
\[
\boxed{A=A^{\mathsf T}}.
\]
\end{solution}

%例题3.1.5
\begin{example}
若 $n$ 阶方阵 $A,B$ 满足
\[
AB-BA=A,
\]
证明:对任意正整数 $k$,都有
\[
\tr(A^k)=0.
\]
\end{example}
\exampleblank

\begin{solution}
先证明
\[
A^kB-BA^k=kA^k.
\]
当 $k=1$ 时即为已知条件.若对 $k$ 成立,则
\[
\begin{aligned}
A^{k+1}B-BA^{k+1}
&=A(A^kB-BA^k)+(AB-BA)A^k\\
&=kA^{k+1}+A^{k+1}\\
&=(k+1)A^{k+1}.
\end{aligned}
\]
故归纳成立.两边取迹得
\[
0=k\tr(A^k),
\]
于是
\[
\boxed{\tr(A^k)=0}.
\]
\end{solution}


\subsection{Kronecker 积(张量积)}

\begin{definition}
设 $A=(a_{ij})$,$B$ 分别为数域 $F$ 上的 $n$ 阶,$m$ 阶方阵,定义
\[
A\otimes B
=
\begin{pmatrix}
a_{11}B&a_{12}B&\cdots&a_{1n}B\\
a_{21}B&a_{22}B&\cdots&a_{2n}B\\
\vdots&\vdots&&\vdots\\
a_{n1}B&a_{n2}B&\cdots&a_{nn}B
\end{pmatrix},
\]
称为 $A$ 与 $B$ 的 Kronecker 积(或张量积).
\end{definition}

\begin{theorem}[\markstar]
若矩阵乘法的尺寸相容,则:
\begin{enumerate}
  \item $(AC)\otimes(BD)=(A\otimes B)(C\otimes D)$;
  \item 若 $A,B$ 可逆,则 $A\otimes B$ 可逆,且
  \[
  (A\otimes B)^{-1}=A^{-1}\otimes B^{-1};
  \]
  \item 若 $A$ 为 $n$ 阶,$B$ 为 $m$ 阶,则
  \[
  |A\otimes B|=|A|^m|B|^n.
  \]
\end{enumerate}
\end{theorem}

\begin{proposition}[\markstar]
设 $A,B$ 分别是 $n,m$ 阶矩阵,$A$ 的特征值为 $\lambda_i$,$B$ 的特征值为 $\mu_j$,则 $A\otimes B$ 的全部特征值为
\[
\lambda_i\mu_j,
\qquad
1\leq i\leq n,
\quad
1\leq j\leq m,
\]
重数按乘积计算.
\end{proposition}

%例题3.1.6
\begin{example}
证明上述 Kronecker 积的三个性质以及特征值命题.
\end{example}
\exampleblank

\begin{solution}
对分块元素直接计算即可得到混合乘积公式
\[
(A\otimes B)(C\otimes D)=(AC)\otimes(BD).
\]
于是若 $A,B$ 可逆,
\[
(A\otimes B)(A^{-1}\otimes B^{-1})=I_n\otimes I_m=I_{nm},
\]
故
\[
(A\otimes B)^{-1}=A^{-1}\otimes B^{-1}.
\]

在代数闭包中,取可逆矩阵 $P,Q$ 使
\[
P^{-1}AP=T_A,
\qquad
Q^{-1}BQ=T_B
\]
均为上三角矩阵.则
\[
(P\otimes Q)^{-1}(A\otimes B)(P\otimes Q)
=T_A\otimes T_B,
\]
后者仍为上三角矩阵,其对角元恰为所有 $\lambda_i\mu_j$.因此
\[
\operatorname{spec}(A\otimes B)
=
\{\lambda_i\mu_j\}_{i,j}.
\]
从而
\[
|A\otimes B|
=
\prod_{i=1}^{n}\prod_{j=1}^{m}\lambda_i\mu_j
=|A|^m|B|^n.
\]
这些恒等式的两端均属于原数域 $F$,故结论成立.
\end{solution}


\subsection{矩阵可交换问题}

若 $AB=BA$,则称方阵 $A,B$ 可交换.

\begin{proposition}[\markstar]
设 $J_n(\lambda)$ 为 $n$ 阶 Jordan 块,则对任意 $n$ 阶方阵 $A$,
\[
AJ_n(\lambda)=J_n(\lambda)A
\]
当且仅当 $A$ 是 $J_n(\lambda)$ 的次数小于 $n$ 的多项式.
\end{proposition}

%例题3.1.7
\begin{example}
求与下列矩阵可交换的全部矩阵:
\begin{enumerate}
  \item $\begin{pmatrix}0&1\\0&0\end{pmatrix}$;
  \item $\begin{pmatrix}0&1&0&0\\0&0&1&0\\0&0&0&1\\0&0&0&0\end{pmatrix}$;
  \item $\begin{pmatrix}3&0&0\\-1&3&0\\0&-1&3\end{pmatrix}$;
  \item $\diag\left(J_2(1),J_2(1)\right)$;
  \item $\diag(1,1,2,2)$.
\end{enumerate}
\end{example}
\exampleblank

\begin{solution}
\leavevmode
\begin{enumerate}
  \item 由 Jordan 块的交换子结构,全部矩阵为
  \[
  \boxed{
  \begin{pmatrix}a&b\\0&a\end{pmatrix}
  }.
  \]
  \item 全部矩阵为四阶上三角 Toeplitz 矩阵:
  \[
  \boxed{
  \begin{pmatrix}
a&b&c&d\\
0&a&b&c\\
0&0&a&b\\
0&0&0&a
\end{pmatrix}
  }.
  \]
  \item 全部矩阵为下三角 Toeplitz 矩阵:
  \[
  \boxed{
  \begin{pmatrix}
a&0&0\\
b&a&0\\
c&b&a
\end{pmatrix}
  }.
  \]
  \item 将未知矩阵按 $2\times2$ 分块写成
  \[
  X=\begin{pmatrix}X_{11}&X_{12}\\X_{21}&X_{22}\end{pmatrix}.
  \]
  每个 $X_{ij}$ 都必须与 $J_2(1)$ 可交换,因此
  \[
  \boxed{
  X_{ij}=\begin{pmatrix}a_{ij}&b_{ij}\\0&a_{ij}\end{pmatrix}
  }.
  \]
  \item 由两个互异特征值对应的块不能相互混合,全部矩阵为
  \[
  \boxed{
  \begin{pmatrix}
  B_1&O\\O&B_2
  \end{pmatrix}
  },
  \]
  其中 $B_1,B_2$ 为任意二阶方阵.
\end{enumerate}
\end{solution}

\begin{theorem}[\marktwostar]
\leavevmode
\begin{enumerate}
  \item 若
  \[
  A=\diag(a_1,a_2,\ldots,a_n),
  \]
  且 $a_i$ 两两不同,则与 $A$ 可交换的矩阵只能是对角矩阵;
  \item 若
  \[
  A=\diag(a_1I_{r_1},a_2I_{r_2},\ldots,a_sI_{r_s}),
  \]
  且 $a_i$ 两两不同,则与 $A$ 可交换的矩阵恰为
  \[
  \diag(B_1,B_2,\ldots,B_s),
  \]
  其中 $B_i$ 为任意 $r_i$ 阶方阵;
  \item 与所有 $n$ 阶可逆矩阵可交换的矩阵恰为数量矩阵 $kI_n$.
\end{enumerate}
\end{theorem}

%例题3.1.8
\begin{example}
证明上述定理.
\end{example}
\exampleblank

\begin{solution}
设 $X=(x_{ij})$.

第一问中,$XA=AX$ 的 $(i,j)$ 元给出
\[
(a_j-a_i)x_{ij}=0.
\]
当 $i\neq j$ 时 $a_i\neq a_j$,故 $x_{ij}=0$,所以 $X$ 为对角矩阵.

第二问按 $r_i\times r_j$ 分块写 $X=(X_{ij})$,则
\[
(a_j-a_i)X_{ij}=O.
\]
因此 $i\neq j$ 时 $X_{ij}=O$,而对角块任意.

第三问中,若 $X$ 与所有可逆矩阵可交换,则对任意矩阵 $Y$,取足够多个标量 $t$ 使 $Y+tI_n$ 可逆,便有
\[
X(Y+tI_n)=(Y+tI_n)X,
\]
从而 $XY=YX$.于是 $X$ 与所有矩阵可交换.特别地,与全部对角矩阵可交换说明 $X$ 为对角矩阵;再与各初等矩阵 $E_{ij}$ 可交换,则所有对角元相等.因此
\[
\boxed{X=kI_n}.
\]
\end{solution}


\newpage

\section{逆矩阵与伴随矩阵}

\begin{theorem}[\markstar]
若 $A,B$ 是同阶方阵,且
\[
AB=I_n,
\]
则
\[
BA=I_n.
\]
\end{theorem}

\begin{proposition}[\markstar]
若
\[
A=\begin{pmatrix}a&b\\c&d\end{pmatrix},
\]
则其伴随矩阵为
\[
A^*=\begin{pmatrix}d&-b\\-c&a\end{pmatrix}.
\]
若 $A$ 可逆,则
\[
A^{-1}
=\dfrac{1}{ad-bc}
\begin{pmatrix}d&-b\\-c&a\end{pmatrix}.
\]
\end{proposition}

%例题3.2.1
\begin{example}
设 $A,B$ 都是 $n$ 阶方阵,且
\[
A+B=AB.
\]
证明:$AB=BA$.
\end{example}
\exampleblank

\begin{solution}
由
\[
A+B=AB
\]
可得
\[
(A-I_n)(B-I_n)=I_n.
\]
根据单侧逆即双侧逆,
\[
(B-I_n)(A-I_n)=I_n.
\]
展开即得
\[
BA-B-A=O.
\]
与已知式比较,
\[
\boxed{AB=BA}.
\]
\end{solution}

%例题3.2.2
\begin{example}
设 $A,B$ 都是 $n$ 阶方阵,且
\[
AB=A-B.
\]
证明:$A$ 可逆的充要条件是 $B$ 可逆.
\end{example}
\exampleblank

\begin{solution}
由已知式
\[
(A+I_n)(B-I_n)=-I_n,
\]
故 $A+I_n$ 与 $B-I_n$ 均可逆.

又
\[
(A+I_n)B=A,
\]
因此
\[
B=(A+I_n)^{-1}A.
\]
所以 $A$ 可逆推出 $B$ 可逆.

另一方面,
\[
A(B-I_n)=-B,
\]
故
\[
A=-B(B-I_n)^{-1}.
\]
所以 $B$ 可逆推出 $A$ 可逆.综上二者等价.
\end{solution}

\begin{theorem}[\markstar]
伴随矩阵具有下列性质:
\begin{enumerate}
  \item $AA^*=A^*A=|A|I_n$;
  \item $(\lambda A)^*=\lambda^{n-1}A^*$;
  \item $|A^*|=|A|^{n-1}$;
  \item $(A^*)^{\mathsf T}=(A^{\mathsf T})^*$;
  \item $(AB)^*=B^*A^*$;
  \item $A$ 可逆时,$A^*=|A|A^{-1}$;
  \item $A$ 可逆时,$(A^*)^{-1}=\dfrac{1}{|A|}A$;
  \item $A$ 可逆时,$(A^*)^{-1}=(A^{-1})^*$;
  \item
  \[
  (A^*)^*
  =
  \begin{cases}
  |A|^{n-2}A,&n>2,\\
  A,&n=2;
  \end{cases}
  \]
  \item
  \[
  \rank(A^*)
  =
  \begin{cases}
  n,&\rank(A)=n,\\
  1,&\rank(A)=n-1,\\
  0,&\rank(A)<n-1.
  \end{cases}
  \]
\end{enumerate}
\end{theorem}

\begin{remark}
处理伴随矩阵恒等式时,常先假设矩阵可逆,利用
\[
A^*=|A|A^{-1}
\]
化简,再通过摄动或多项式恒等原理推广到一般矩阵.这是高等代数中常用的``摄动法''思路.
\end{remark}

%例题3.2.3
\begin{example}
设 $A,B$ 为 $n$ 阶方阵,证明:
\[
(AB)^*=B^*A^*.
\]
\end{example}
\exampleblank

\begin{solution}
若 $A,B$ 均可逆,则
\[
(AB)^*
=|AB|(AB)^{-1}
=|A||B|B^{-1}A^{-1}
=B^*A^*.
\]

一般情形下,将 $A,B$ 分别替换为 $A+tI_n,B+sI_n$.除有限多个 $t,s$ 外它们均可逆,于是上述恒等式在无穷多个参数点成立.两边各元素都是 $t,s$ 的多项式,因此恒等式恒成立.令 $t=s=0$ 即得结论.
\end{solution}

%例题3.2.4
\begin{example}
设 $A$ 为 $n$ 阶方阵,证明:
\[
|A^*|=|A|^{n-1}.
\]
\end{example}
\exampleblank

\begin{solution}
若 $A$ 可逆,则
\[
A^*=|A|A^{-1},
\]
故
\[
|A^*|
=|A|^n|A^{-1}|
=|A|^{n-1}.
\]
一般情形可由摄动法或注意到两端均为 $A$ 的元素的多项式而推广得到.
\end{solution}

%例题3.2.5
\begin{example}
设 $A,B,C,D$ 都是 $n$ 阶方阵,且 $AC=CA$.证明:
\[
\begin{vmatrix}
A&B\\C&D
\end{vmatrix}
=|AD-CB|.
\]
\end{example}
\exampleblank

\begin{solution}
先设 $A$ 可逆.由 Schur 补公式,
\[
\begin{vmatrix}A&B\\C&D\end{vmatrix}
=|A|\,|D-CA^{-1}B|.
\]
由于 $AC=CA$,故 $A^{-1}C=CA^{-1}$,于是
\[
A(D-CA^{-1}B)=AD-CB.
\]
因此
\[
|A|\,|D-CA^{-1}B|=|AD-CB|.
\]

若 $A$ 不可逆,考虑 $A+tI_n$.它仍与 $C$ 可交换,且除有限多个 $t$ 外可逆.上式对无穷多个 $t$ 成立,故由多项式恒等原理对 $t=0$ 也成立.
\end{solution}

%例题3.2.6
\begin{example}
设 $A,B,C,D$ 都是 $n$ 阶方阵,且
\[
AB^{\mathsf T}=BA^{\mathsf T}.
\]
证明:
\[
\begin{vmatrix}
A&B\\C&D
\end{vmatrix}
=|AD^{\mathsf T}-BC^{\mathsf T}|.
\]
\end{example}
\exampleblank

\begin{solution}
先设 $A$ 可逆.由条件
\[
AB^{\mathsf T}=BA^{\mathsf T}
\]
得
\[
B^{\mathsf T}=A^{-1}BA^{\mathsf T}.
\]
故
\[
\begin{aligned}
|AD^{\mathsf T}-BC^{\mathsf T}|
&=|DA^{\mathsf T}-CB^{\mathsf T}|\\
&=|(D-CA^{-1}B)A^{\mathsf T}|\\
&=|A|\,|D-CA^{-1}B|\\
&=
\begin{vmatrix}A&B\\C&D\end{vmatrix}.
\end{aligned}
\]
一般奇异情形可由该块行列式恒等式的多项式延拓得到.
\end{solution}

%例题3.2.7
\begin{example}
设 $A,B,C,D$ 都是 $n$ 阶矩阵,若
\[
M=\begin{pmatrix}A&B\\C&D\end{pmatrix}
\]
的秩为 $n$,证明:
\[
\begin{vmatrix}|A|&|B|\\|C|&|D|\end{vmatrix}=0.
\]
且若 $A$ 可逆,则
\[
D=CA^{-1}B.
\]
\end{example}
\exampleblank

\begin{solution}
若 $A$ 可逆,对 $M$ 作可逆块初等变换:
\[
\begin{pmatrix}I&O\\-CA^{-1}&I\end{pmatrix}
\begin{pmatrix}A&B\\C&D\end{pmatrix}
=
\begin{pmatrix}A&B\\O&D-CA^{-1}B\end{pmatrix}.
\]
因此
\[
\rank(M)=n+\rank(D-CA^{-1}B).
\]
已知 $\rank(M)=n$,故
\[
\boxed{D=CA^{-1}B}.
\]
进而
\[
|A||D|=|A|\,|CA^{-1}B|=|C||B|.
\]

若 $A$ 不可逆,则 $|A|=0$.若同时 $|B|\neq0$ 且 $|C|\neq0$,利用 $B$ 可逆作同样的块消元可得
\[
C=DB^{-1}A,
\]
右端奇异,与 $C$ 可逆矛盾.因此 $|B||C|=0$,仍有
\[
|A||D|-|B||C|=0.
\]
\end{solution}

%例题3.2.8
\begin{example}
设
\[
J=J_n(1)=
\begin{pmatrix}
1&1&&&\\
&1&1&&\\
&&\ddots&\ddots&\\
&&&1&1\\
&&&&1
\end{pmatrix}.
\]
求解矩阵方程
\[
3X=XJ+JX.
\]
\end{example}
\exampleblank

\begin{solution}
写
\[
J=I_n+N,
\]
其中 $N$ 的超对角线上全为 $1$.方程化为
\[
X=XN+NX.
\]
设 $X=(x_{ij})$,约定 $x_{i0}=0,x_{n+1,j}=0$,则
\[
x_{ij}=x_{i,j-1}+x_{i+1,j}.
\]
从 $i=n,j=1$ 开始有 $x_{n1}=0$,继而递推得到第 $n$ 行全为零;再向上逐行递推,最终所有元素均为零.因此
\[
\boxed{X=O}.
\]
\end{solution}

%例题3.2.9
\begin{example}
设 $A$ 为 $n$ 阶可逆矩阵,$\alpha,\beta$ 均为 $n$ 维列向量,且
\[
1+\beta^{\mathsf T}A^{-1}\alpha\neq0.
\]
\begin{enumerate}
  \item 证明
  \[
  P=\begin{pmatrix}A&\alpha\\-\beta^{\mathsf T}&1\end{pmatrix}
  \]
  可逆,并求 $P^{-1}$;
  \item 证明
  \[
  Q=A+\alpha\beta^{\mathsf T}
  \]
  可逆,并求 $Q^{-1}$.
\end{enumerate}
\end{example}
\exampleblank

\begin{solution}
令
\[
s=1+\beta^{\mathsf T}A^{-1}\alpha.
\]
由于 $s\neq0$,由 Schur 补公式,$P$ 可逆,且
\[
\boxed{
P^{-1}
=
\begin{pmatrix}
A^{-1}-\dfrac{A^{-1}\alpha\beta^{\mathsf T}A^{-1}}{s}
&-\dfrac{A^{-1}\alpha}{s}\\[8pt]
\dfrac{\beta^{\mathsf T}A^{-1}}{s}
&\dfrac{1}{s}
\end{pmatrix}
}.
\]

第二问为 Sherman--Morrison 公式:
\[
\boxed{
Q^{-1}
=A^{-1}
-\dfrac{A^{-1}\alpha\beta^{\mathsf T}A^{-1}}
{1+\beta^{\mathsf T}A^{-1}\alpha}
}.
\]
直接相乘即可验证.
\end{solution}

%例题3.2.10
\begin{example}
设 $A,B,C$ 都是 $n$ 阶方阵,
\[
Q=
\begin{pmatrix}
A&A\\
C-B&C
\end{pmatrix}.
\]
\begin{enumerate}
  \item 证明:$Q$ 可逆的充要条件是 $AB$ 可逆;
  \item 当 $Q$ 可逆时,求 $Q^{-1}$.
\end{enumerate}
\end{example}
\exampleblank

\begin{solution}
作块列变换``第二块列减第一块列'',即右乘
\[
R=\begin{pmatrix}I_n&-I_n\\O&I_n\end{pmatrix},
\]
得到
\[
QR
=
\begin{pmatrix}
A&O\\
C-B&B
\end{pmatrix}.
\]
因此 $Q$ 可逆当且仅当 $A,B$ 都可逆,也即当且仅当 $AB$ 可逆.

此时
\[
(QR)^{-1}
=
\begin{pmatrix}
A^{-1}&O\\
-B^{-1}(C-B)A^{-1}&B^{-1}
\end{pmatrix}.
\]
因为 $Q^{-1}=R(QR)^{-1}$,故
\[
\boxed{
Q^{-1}
=
\begin{pmatrix}
B^{-1}CA^{-1}&-B^{-1}\\
B^{-1}(B-C)A^{-1}&B^{-1}
\end{pmatrix}
}.
\]
\end{solution}

%例题3.2.11
\begin{example}
设 $A$ 是 $n$ 阶方阵($n\geq2$),且存在 $n$ 阶方阵 $B$,使得
\[
A^*B-BA^*=A^*.
\]
证明:
\begin{enumerate}
  \item $|A|=0$;
  \item $(A^*)^2=O$.
\end{enumerate}
\end{example}
\exampleblank

\begin{solution}
令
\[
C=A^*.
\]
则
\[
CB-BC=C.
\]
由例题 3.1.5,
\[
\tr(C^k)=0,
\qquad k=1,2,\ldots.
\]
若 $C$ 可逆,则其特征值全非零;但 Newton 恒等式结合 $\tr(C^k)=0\,(1\leq k\leq n)$ 会推出 $C$ 的特征多项式为 $x^n$,矛盾.因此 $|C|=0$.又
\[
|C|=|A^*|=|A|^{n-1},
\]
故
\[
\boxed{|A|=0}.
\]
于是 $\rank(A)\leq n-1$,由伴随矩阵秩的性质,$\rank(C)\leq1$.又由原方程取迹得
\[
\tr(C)=0.
\]
若 $C=uv^{\mathsf T}$ 为秩一矩阵,则
\[
C^2=(v^{\mathsf T}u)C=\tr(C)C=O.
\]
若 $C=O$ 更显然.因此
\[
\boxed{(A^*)^2=O}.
\]
\end{solution}

%例题3.2.12
\begin{example}[\markstar]
设 $A,B$ 是 $n$ 阶矩阵,$I_n+AB$ 可逆.证明:$I_n+BA$ 也可逆.
\end{example}
\exampleblank

\begin{solution}
若
\[
(I_n+BA)x=0,
\]
则
\[
(I_n+AB)Ax=A(I_n+BA)x=0.
\]
由于 $I_n+AB$ 可逆,故 $Ax=0$.代回原式得到
\[
x=-BAx=0.
\]
因此 $I_n+BA$ 的核只有零向量,作为方阵它必可逆.
\end{solution}


\newpage

\section{等价标准型}

\begin{theorem}[\marktwostar]
设 $A_{m\times n}$ 满足 $\rank(A)=r$,则存在 $m$ 阶可逆矩阵 $P$ 与 $n$ 阶可逆矩阵 $Q$,使
\[
A=P
\begin{pmatrix}
I_r&O\\O&O
\end{pmatrix}
Q.
\]
\end{theorem}

\begin{remark}
该标准型常用于处理``只知道矩阵的秩而没有更多结构''的问题.
\end{remark}

%例题3.3.1
\begin{example}
设 $A$ 是 $n$ 阶方阵.证明:存在 $n$ 阶可逆矩阵 $B$ 以及 $n$ 阶幂等矩阵 $C$(即 $C^2=C$),使得
\[
A=BC.
\]
\end{example}
\exampleblank

\begin{solution}
设 $\rank(A)=r$.由等价标准型,存在可逆矩阵 $P,Q$ 使
\[
A=P E_r Q,
\qquad
E_r=\begin{pmatrix}I_r&O\\O&O\end{pmatrix}.
\]
令
\[
B=PQ,
\qquad
C=Q^{-1}E_rQ.
\]
则 $B$ 可逆,且
\[
C^2=Q^{-1}E_r^2Q=C,
\]
同时
\[
BC=PQ\,Q^{-1}E_rQ=PE_rQ=A.
\]
\end{solution}

%例题3.3.2
\begin{example}[\markstar]
设矩阵 $A_{m\times n}$ 行满秩,即 $\rank(A)=m$.证明:
\begin{enumerate}
  \item 存在可逆矩阵 $Q$,使
  \[
  A=(I_m,O)Q;
  \]
  \item 存在矩阵 $B_{n\times m}$,使
  \[
  AB=I_m.
  \]
\end{enumerate}
\end{example}
\exampleblank

\begin{solution}
由等价标准型,
\[
A=P(I_m,O)Q_0
\]
其中 $P,Q_0$ 可逆.令
\[
Q=\begin{pmatrix}P&O\\O&I_{n-m}\end{pmatrix}Q_0,
\]
则 $Q$ 可逆,且
\[
A=(I_m,O)Q.
\]
再令
\[
B=Q^{-1}\begin{pmatrix}I_m\\O\end{pmatrix},
\]
则
\[
AB=(I_m,O)QQ^{-1}\begin{pmatrix}I_m\\O\end{pmatrix}=I_m.
\]
\end{solution}

%例题3.3.3
\begin{example}[\markstar]
\begin{enumerate}
  \item 证明:若 $A_{m\times n}$ 且 $\rank(A)=r$,则存在列满秩矩阵 $B_{m\times r}$ 和行满秩矩阵 $C_{r\times n}$,使
  \[
  A=BC.
  \]
  这称为 $A$ 的满秩分解;
  \item 求矩阵
  \[
  A=
  \begin{pmatrix}
  1&2&0&1&1&10\\
  3&6&1&4&2&36\\
  2&4&0&2&2&27\\
  6&12&1&7&5&73
  \end{pmatrix}
  \]
  的一个满秩分解;
  \item 若 $A=BC$ 是满秩分解,证明 $Ax=0$ 与 $Cx=0$ 同解.
\end{enumerate}
\end{example}
\exampleblank

\begin{solution}
\leavevmode
\begin{enumerate}
  \item 从 $A$ 中取 $r$ 个线性无关列组成 $B$.由于 $B$ 的列构成 $A$ 的列空间的一组基,每一列 $A_j$ 都可唯一写成 $B$ 的列的线性组合.把这些坐标作为 $C$ 的各列,即有
  \[
  A=BC.
  \]
  因为所选的 $r$ 个基列对应的坐标列构成 $I_r$,故 $C$ 行满秩.

  \item 该矩阵的秩为 $3$.取第 $1,3,6$ 列作为基列,令
  \[
  B=
  \begin{pmatrix}
  1&0&10\\
  3&1&36\\
  2&0&27\\
  6&1&73
  \end{pmatrix},
  \qquad
  C=
  \begin{pmatrix}
  1&2&0&1&1&0\\
  0&0&1&1&-1&0\\
  0&0&0&0&0&1
  \end{pmatrix}.
  \]
  直接相乘得到
  \[
  \boxed{A=BC}.
  \]

  \item 因 $B$ 列满秩,线性映射 $y\mapsto By$ 的核为零.因此
  \[
  Ax=0
  \iff BCx=0
  \iff Cx=0.
  \]
\end{enumerate}
\end{solution}

%例题3.3.4
\begin{example}
设 $A$ 是 $n$ 阶方阵且 $|A|=0$.证明:存在非零 $n$ 阶方阵 $B$,使
\[
AB=BA=O.
\]
\end{example}
\exampleblank

\begin{solution}
因为 $A$ 奇异,存在非零列向量 $u,v$ 使
\[
Au=0,
\qquad
A^{\mathsf T}v=0.
\]
令
\[
B=uv^{\mathsf T}.
\]
则 $B\neq O$,且
\[
AB=(Au)v^{\mathsf T}=O,
\]
同时
\[
BA=u(v^{\mathsf T}A)=u(A^{\mathsf T}v)^{\mathsf T}=O.
\]
\end{solution}

%例题3.3.5
\begin{example}
证明:任意秩为 $r>0$ 的矩阵都可以分解为 $r$ 个秩为 $1$ 的矩阵之和.
\end{example}
\exampleblank

\begin{solution}
取满秩分解
\[
A=BC,
\]
其中
\[
B=(b_1,b_2,\ldots,b_r),
\qquad
C=
\begin{pmatrix}
c_1^{\mathsf T}\\
c_2^{\mathsf T}\\
\vdots\\
c_r^{\mathsf T}
\end{pmatrix}.
\]
则
\[
A
=\sum\limits_{i=1}^{r}b_ic_i^{\mathsf T}.
\]
由于 $B$ 列满秩,$C$ 行满秩,每个 $b_i,c_i$ 均非零,所以每个 $b_ic_i^{\mathsf T}$ 的秩都为 $1$.
\end{solution}

%例题3.3.6
\begin{example}
若 $A^2=O$,证明存在可逆矩阵 $P$,使
\[
P^{-1}AP
=
\begin{pmatrix}
O&I_r&O\\
O&O&O\\
O&O&O
\end{pmatrix},
\qquad r=\rank(A).
\]
\end{example}
\exampleblank

\begin{solution}
因为 $A^2=O$,有
\[
\operatorname{Im}A\subseteq\ker A.
\]
取
\[
v_1,\ldots,v_r
\]
为 $\operatorname{Im}A$ 的一组基.对每个 $v_i$ 取 $u_i$ 使
\[
Au_i=v_i.
\]
再把 $v_1,\ldots,v_r$ 扩充为 $\ker A$ 的一组基
\[
v_1,\ldots,v_r,w_1,\ldots,w_{n-2r}.
\]
则
\[
v_1,\ldots,v_r,u_1,\ldots,u_r,w_1,\ldots,w_{n-2r}
\]
构成全空间的一组基.在这组基下,$A$ 的矩阵恰为
\[
\boxed{
\begin{pmatrix}
O&I_r&O\\
O&O&O\\
O&O&O
\end{pmatrix}
}.
\]
\end{solution}

%例题3.3.7
\begin{example}
已知 $n$ 阶矩阵 $A$ 的 $n$ 个顺序主子式均非零.证明:存在 $n$ 阶下三角矩阵 $B$,使 $BA$ 为上三角矩阵.
\end{example}
\exampleblank

\begin{solution}
条件``所有顺序主子式非零''恰好保证 Gaussian 消元过程中不需要交换行,因此 $A$ 存在分解
\[
A=LU,
\]
其中 $L$ 为可逆下三角矩阵,$U$ 为上三角矩阵.令
\[
B=L^{-1},
\]
则 $B$ 仍为下三角矩阵,并且
\[
\boxed{BA=U}
\]
为上三角矩阵.
\end{solution}

%例题3.3.8
\begin{example}
设 $A$ 为 $n$ 阶实方阵,$A$ 的所有顺序主子式都大于零,且所有非主对角元均小于零.证明:$A^{-1}$ 的所有元素都大于零.
\end{example}
\exampleblank

\begin{solution}
对 $n$ 作归纳.$n=1$ 显然.设结论对 $n-1$ 阶成立,将 $A$ 分块为
\[
A=
\begin{pmatrix}
A_{n-1}&b\\
c^{\mathsf T}&a
\end{pmatrix}.
\]
由条件,$A_{n-1}$ 同样满足归纳假设,因此 $A_{n-1}^{-1}$ 的所有元素均为正,而 $b,c$ 的所有元素均为负.

Schur 补
\[
s=a-c^{\mathsf T}A_{n-1}^{-1}b
\]
满足
\[
s=\dfrac{|A|}{|A_{n-1}|}>0.
\]
块逆公式给出
\[
A^{-1}
=
\begin{pmatrix}
A_{n-1}^{-1}+A_{n-1}^{-1}bs^{-1}c^{\mathsf T}A_{n-1}^{-1}
&-A_{n-1}^{-1}bs^{-1}\\
-s^{-1}c^{\mathsf T}A_{n-1}^{-1}
&s^{-1}
\end{pmatrix}.
\]
其中 $A_{n-1}^{-1}b$ 与 $c^{\mathsf T}A_{n-1}^{-1}$ 的元素均为负,故上式四个块的元素全部严格为正.因此 $A^{-1}$ 的所有元素均大于零.
\end{solution}


\newpage

\section{矩阵方程}

%例题3.4.1
\begin{example}
求解
\[
X
\begin{pmatrix}
1&1&-2\\
0&1&3\\
1&0&0
\end{pmatrix}
=
\begin{pmatrix}
1&-1&1\\
0&1&0\\
1&0&2
\end{pmatrix}.
\]
\end{example}
\exampleblank

\begin{solution}
右乘系数矩阵的逆,得到
\[
\boxed{
X=
\begin{pmatrix}
-\dfrac45&-\dfrac15&\dfrac95\\[4pt]
\dfrac35&\dfrac25&-\dfrac35\\[4pt]
-\dfrac25&\dfrac25&\dfrac75
\end{pmatrix}
}.
\]
\end{solution}

%例题3.4.2
\begin{example}
已知
\[
A=
\begin{pmatrix}
-2&1&1\\
1&-2&1\\
1&1&-2
\end{pmatrix},
\qquad
B=
\begin{pmatrix}
1&b\\
2&a\\
a&2
\end{pmatrix},
\]
且矩阵方程 $AX=B$ 有解.求 $a,b,X$.
\end{example}
\exampleblank

\begin{solution}
由于
\[
(1,1,1)A=0,
\]
方程有解必须满足 $B$ 的每一列与左零空间正交,即
\[
1+2+a=0,
\qquad
b+a+2=0.
\]
故
\[
\boxed{a=-3,\qquad b=1}.
\]
此时通解为
\[
\boxed{
X=
\begin{pmatrix}
s-\dfrac43&t+\dfrac13\\[4pt]
s-\dfrac53&t+\dfrac53\\[4pt]
s&t
\end{pmatrix},
\qquad s,t\in F
}.
\]
\end{solution}

%例题3.4.3
\begin{example}
解下列矩阵方程:
\begin{enumerate}
  \item
  \[
  X
  \begin{pmatrix}
  2&0&0\\
  0&2&5\\
  0&3&8
  \end{pmatrix}
  =
  \begin{pmatrix}
  1&-1&1\\
  2&-3&1\\
  3&-4&1
  \end{pmatrix};
  \]
  \item
  \[
  \begin{pmatrix}4&6\\6&9\end{pmatrix}X
  =
  \begin{pmatrix}1&1\\1&1\end{pmatrix};
  \]
  \item
  \[
  \begin{pmatrix}
  2&-3&1\\
  4&-5&2\\
  5&-7&3
  \end{pmatrix}
  X
  \begin{pmatrix}
  9&7&6\\
  1&1&2\\
  1&1&1
  \end{pmatrix}
  =
  \begin{pmatrix}
  2&0&-2\\
  18&12&9\\
  23&15&11
  \end{pmatrix};
  \]
  \item 设
  \[
  U=
  \begin{pmatrix}
  1&1&1&\cdots&1\\
  0&1&1&\cdots&1\\
  0&0&1&\cdots&1\\
  \vdots&\vdots&\vdots&\ddots&\vdots\\
  0&0&0&\cdots&1
  \end{pmatrix},
  \]
  求解
  \[
  UX=
  \begin{pmatrix}
  1&2&3&\cdots&n\\
  0&1&2&\cdots&n-1\\
  0&0&1&\cdots&n-2\\
  \vdots&\vdots&\vdots&\ddots&\vdots\\
  0&0&0&\cdots&1
  \end{pmatrix}.
  \]
\end{enumerate}
\end{example}
\exampleblank

\begin{solution}
\leavevmode
\begin{enumerate}
  \item 系数矩阵可逆,直接右乘其逆得
  \[
  \boxed{
  X=
  \begin{pmatrix}
  \dfrac12&-11&7\\
  1&-27&17\\
  \dfrac32&-35&22
  \end{pmatrix}
  }.
  \]
  \item 左侧系数矩阵的列空间由 $(2,3)^{\mathsf T}$ 张成,而右端每一列为 $(1,1)^{\mathsf T}$,不属于该列空间,故
  \[
  \boxed{\text{无解}}.
  \]
  \item 左右两个系数矩阵均可逆,因此
  \[
  \boxed{
  X=
  \begin{pmatrix}
  1&1&1\\
  1&2&3\\
  2&3&1
  \end{pmatrix}
  }.
  \]
  \item 右端矩阵恰为 $U^2$,故
  \[
  \boxed{X=U}.
  \]
\end{enumerate}
\end{solution}

\begin{proposition}[\markstar]
设 $A,B$ 分别为数域 $F$ 上的 $n,m$ 阶矩阵.若 $A,B$ 有 $r$ 个互异公共特征值,则矩阵方程
\[
AX-XB=O
\]
存在秩为 $r$ 的矩阵解.
\end{proposition}

\begin{proposition}[\markstar]
设 $A,B$ 分别为复数域上的 $n,m$ 阶矩阵.若矩阵方程
\[
AX-XB=O
\]
存在秩为 $r$ 的矩阵解,则 $A,B$ 至少有 $r$ 个公共特征值,重根按重数计算.
\end{proposition}

\begin{proposition}[\marktwostar]
设 $A,B$ 分别为复数域上的 $n,m$ 阶矩阵,则矩阵方程
\[
AX=XB
\]
只有零解的充要条件是 $A,B$ 没有公共特征值.
\end{proposition}

%例题3.4.4
\begin{example}
证明上述三个命题.
\end{example}
\exampleblank

\begin{solution}
对第一个命题,设 $\lambda_1,\ldots,\lambda_r$ 是互异公共特征值.取
\[
A\alpha_i=\lambda_i\alpha_i,
\qquad
\beta_i^{\mathsf T}B=\lambda_i\beta_i^{\mathsf T},
\]
其中 $\alpha_i$ 为 $A$ 的右特征向量,$\beta_i^{\mathsf T}$ 为 $B$ 的左特征向量.令
\[
X=\sum\limits_{i=1}^{r}\alpha_i\beta_i^{\mathsf T}.
\]
则 $AX=XB$.由于互异特征值对应的左右特征向量组分别线性无关,故 $\rank(X)=r$.

对第二个命题,把 $X$ 看作线性映射
\[
X:\mathbb{C}^m\longrightarrow\mathbb{C}^n.
\]
由 $AX=XB$,$\operatorname{Im}X$ 是 $A$-不变子空间,而 $\ker X$ 是 $B$-不变子空间.诱导映射
\[
\bar X:\mathbb{C}^m/\ker X\longrightarrow\operatorname{Im}X
\]
是一个 $r$ 维空间之间的同构,并且 intertwine 两边的诱导算子.因此 $A|_{\operatorname{Im}X}$ 与 $B$ 在商空间上的诱导算子相似,具有相同的 $r$ 次特征多项式.该多项式分别整除 $A,B$ 的特征多项式,故二者至少有 $r$ 个公共特征值(计重数).

第三个命题立即由前两个命题得到:若有公共特征值,则第一命题给出非零解;若存在非零解,则其秩至少为 $1$,第二命题推出存在公共特征值.
\end{solution}

%例题3.4.5
\begin{example}
设
\[
A=J_m(a)^{\mathsf T},
\qquad
B=J_n(b)^{\mathsf T},
\]
即二者主对角线分别为 $a,b$,次对角线为 $1$.若
\[
AX=XB,
\]
证明:当 $a\neq b$ 时,$X=O$.
\end{example}
\exampleblank

\begin{solution}
$A$ 的全部特征值均为 $a$,$B$ 的全部特征值均为 $b$.当 $a\neq b$ 时,两者没有公共特征值.由上一命题,矩阵方程
\[
AX=XB
\]
只有零解.因此
\[
\boxed{X=O}.
\]
\end{solution}

%例题3.4.6
\begin{example}
设 $J=J_n(1)$($n\geq2$),求解
\[
3X=XJ+JX.
\]
\end{example}
\exampleblank

\begin{solution}
与例题 3.2.8 相同,唯一解为
\[
\boxed{X=O}.
\]
\end{solution}

%例题3.4.7
\begin{example}
设 $A,B$ 为 $n$ 阶实矩阵,并且 $A,B$ 的全部特征值都大于零.证明:若
\[
A^2=B^2,
\]
则 $A=B$.
\end{example}
\exampleblank

\begin{solution}
令
\[
X=A-B.
\]
则
\[
AX+XB
=A(A-B)+(A-B)B
=A^2-B^2=O.
\]
即
\[
AX=X(-B).
\]
矩阵 $A$ 的特征值全部为正,而 $-B$ 的特征值全部为负,因此二者没有公共特征值.由上一命题,$X=O$.所以
\[
\boxed{A=B}.
\]
\end{solution}

%例题3.4.8
\begin{example}
设 $A_{m\times n}$,求矩阵 $X$ 使
\[
A^{\mathsf T}X=X^{\mathsf T}A.
\]
\end{example}
\exampleblank

\begin{solution}
该方程要求 $A^{\mathsf T}X$ 为对称矩阵,因此解一般不唯一.最直接的一族解是
\[
\boxed{X=AS},
\]
其中 $S$ 满足
\[
A^{\mathsf T}AS=S^{\mathsf T}A^{\mathsf T}A.
\]
特别地,取任意标量 $c$,都有
\[
\boxed{X=cA}
\]
满足条件;尤其 $X=A$ 是一个非零自然解(当 $A\neq O$).

若 $A$ 列满秩,则还可以写出全部解:令
\[
G=A^{\mathsf T}A
\]
可逆,则任取对称矩阵 $H$ 以及满足 $A^{\mathsf T}Z=O$ 的矩阵 $Z$,
\[
\boxed{X=AG^{-1}H+Z}
\]
给出全部解.
\end{solution}

%例题3.4.9
\begin{example}[\markstar\ Roth 定理]
设 $A_{m\times n}$,$B_{p\times q}$,$C_{m\times q}$.证明矩阵方程
\[
AX-YB=C
\]
有解 $X,Y$ 的充要条件是
\[
\rank
\begin{pmatrix}A&O\\O&B\end{pmatrix}
=
\rank
\begin{pmatrix}A&C\\O&B\end{pmatrix}.
\]
\end{example}
\exampleblank

\begin{solution}
若存在 $X,Y$ 使 $AX-YB=C$,则
\[
\begin{pmatrix}I_m&-Y\\O&I_p\end{pmatrix}
\begin{pmatrix}A&O\\O&B\end{pmatrix}
\begin{pmatrix}I_n&X\\O&I_q\end{pmatrix}
=
\begin{pmatrix}A&C\\O&B\end{pmatrix}.
\]
左右两个因子均可逆,因此两块矩阵秩相等.

反过来,分别用可逆行,列变换把 $A,B$ 化为秩标准型
\[
A\sim
\begin{pmatrix}I_r&O\\O&O\end{pmatrix},
\qquad
B\sim
\begin{pmatrix}I_s&O\\O&O\end{pmatrix}.
\]
对应地把 $C$ 分块.此时右侧块矩阵与对角块矩阵秩相等,当且仅当 $C$ 中同时落在 $A$ 的零行部分与 $B$ 的零列部分的右下角块为零.该条件恰好保证其余三个块分别可以写成 $AX$ 或 $YB$ 的相应块,因此可以构造 $X,Y$ 使
\[
AX-YB=C.
\]
故秩条件也是充分的.
\end{solution}


\newpage

\section{广义逆}

\subsection{加号广义逆}

\begin{definition}
设 $A$ 为 $m\times n$ 实矩阵.满足
\begin{enumerate}
  \item $AA^+A=A$;
  \item $A^+AA^+=A^+$;
  \item $(AA^+)^{\mathsf T}=AA^+$;
  \item $(A^+A)^{\mathsf T}=A^+A$
\end{enumerate}
的矩阵 $A^+$ 称为 $A$ 的 Moore--Penrose 广义逆.

在复数情形中,上述转置应改为共轭转置 $\mathrm H$.
\end{definition}

\begin{theorem}[\markstar]
任意实矩阵 $A$ 的 Moore--Penrose 广义逆存在且唯一.若
\[
A=CR
\]
是 $A$ 的满秩分解,其中 $C$ 列满秩,$R$ 行满秩,则
\[
\boxed{
A^+
=R^{\mathsf T}(RR^{\mathsf T})^{-1}(C^{\mathsf T}C)^{-1}C^{\mathsf T}
}.
\]
复数情形相应把转置改为共轭转置.
\end{theorem}

%例题3.5.1
\begin{example}
证明上述 Moore--Penrose 广义逆公式以及唯一性.
\end{example}
\exampleblank

\begin{solution}
记
\[
R^+=R^{\mathsf T}(RR^{\mathsf T})^{-1},
\qquad
C^+=(C^{\mathsf T}C)^{-1}C^{\mathsf T}.
\]
则
\[
RR^+=I_r,
\qquad
C^+C=I_r.
\]
令
\[
X=R^+C^+.
\]
于是
\[
AXA=CRR^+C^+CR=CR=A,
\]
以及
\[
XAX=R^+C^+CRR^+C^+=R^+C^+=X.
\]
同时
\[
AX=CC^+,
\qquad
XA=R^+R
\]
分别是到 $\operatorname{col}(C)$ 与 $\operatorname{row}(R)$ 的正交投影,因此都为对称矩阵.故 $X$ 满足四个 Penrose 方程,即 $X=A^+$.

若 $X,Y$ 都满足四个 Penrose 方程,则 $AX,AY$ 都是到 $\operatorname{col}(A)$ 的正交投影,故 $AX=AY$;同理 $XA=YA$.于是
\[
X=XAX=XAY=YAY=Y.
\]
故 Moore--Penrose 广义逆唯一.
\end{solution}

\begin{proposition}[\markstar]
\leavevmode
\begin{enumerate}
  \item $(A^+)^+=A$;
  \item
  \[
  \rank(A)=\rank(A^+)=\rank(A^+A)=\rank(AA^+).
  \]
\end{enumerate}
\end{proposition}

%例题3.5.2
\begin{example}
证明上述两个性质.
\end{example}
\exampleblank

\begin{solution}
四个 Penrose 方程在交换 $A$ 与 $A^+$ 后保持同样的形式,因此 $A$ 是 $A^+$ 的 Moore--Penrose 广义逆.由唯一性,
\[
\boxed{(A^+)^+=A}.
\]
由
\[
AA^+A=A
\]
有
\[
\rank(A)
\leq\rank(AA^+)
\leq\rank(A),
\]
故 $\rank(AA^+)=\rank(A)$.类似地,由 $A^+AA^+=A^+$ 得
\[
\rank(A^+A)=\rank(A^+).
\]
又
\[
\rank(AA^+)\leq\rank(A^+),
\qquad
\rank(A^+A)\leq\rank(A),
\]
因此
\[
\boxed{
\rank(A)=\rank(A^+)=\rank(A^+A)=\rank(AA^+)
}.
\]
\end{solution}

%例题3.5.3
\begin{example}
求下列矩阵的 Moore--Penrose 广义逆:
\begin{enumerate}
  \item $O_{1\times1}$;
  \item $\diag(\lambda_1,\ldots,\lambda_r,0,\ldots,0)$,其中 $\lambda_i\neq0$;
  \item $\begin{pmatrix}I_r&O\\O&O\end{pmatrix}$;
  \item $\begin{pmatrix}I_r\\O\end{pmatrix}$;
  \item $(I_r,O)$.
\end{enumerate}
\end{example}
\exampleblank

\begin{solution}
分别有
\begin{enumerate}
  \item $\boxed{A^+=0}$;
  \item
  \[
  \boxed{
  A^+=\diag(\lambda_1^{-1},\ldots,\lambda_r^{-1},0,\ldots,0)
  };
  \]
  \item
  \[
  \boxed{
  A^+=\begin{pmatrix}I_r&O\\O&O\end{pmatrix}
  };
  \]
  \item
  \[
  \boxed{A^+=(I_r,O)};
  \]
  \item
  \[
  \boxed{A^+=\begin{pmatrix}I_r\\O\end{pmatrix}}.
  \]
\end{enumerate}
\end{solution}


\subsection{减号广义逆}

\begin{definition}
满足
\[
AA^-A=A
\]
的矩阵 $A^-$ 称为 $A$ 的减号广义逆(或一般广义逆).
\end{definition}

\begin{theorem}[\markstar]
设复矩阵 $A$ 满足
\[
PAQ=
\begin{pmatrix}I_r&O\\O&O\end{pmatrix}
\]
其中 $P,Q$ 可逆,则 $A$ 的所有减号广义逆为
\[
\boxed{
A^-
=Q
\begin{pmatrix}
I_r&Y_2\\
Y_3&Y_4
\end{pmatrix}
P
},
\]
其中 $Y_2,Y_3,Y_4$ 的元素任意.
\end{theorem}

\begin{remark}
只满足 Penrose 条件 (1),(3) 的广义逆常记为 $A_l^-$,称为最小二乘广义逆;只满足条件 (1),(4) 的常记为 $A_m^-$,称为极小范数广义逆.
\end{remark}

%例题3.5.4
\begin{example}
求下列矩阵的一个减号广义逆 $A^-$ 和 Moore--Penrose 广义逆 $A^+$:
\begin{enumerate}
  \item $(a,b,c)$;
  \item $(a,b,c)^{\mathsf T}$;
  \item $\diag(a,b,0)$;
  \item $\begin{pmatrix}a&b\\2a&2b\end{pmatrix}$;
  \item $\begin{pmatrix}a&b&c\\e&f&g\\0&0&0\end{pmatrix}$;
  \item $\begin{pmatrix}a&e&0\\b&f&0\\c&g&0\end{pmatrix}$;
  \item $\begin{pmatrix}a&b&c\\2a&2b&2c\\3a&3b&3c\end{pmatrix}$;
  \item $J_n(0)$;
  \item $E_{1n}$.
\end{enumerate}
\end{example}
\exampleblank

\begin{solution}
以下给出一个减号广义逆;由于 $A^+$ 本身总满足 $AA^+A=A$,可以直接取
\[
A^-=A^+.
\]
因此只需写出 $A^+$.参数退化导致秩下降时,应按实际秩重新计算.

\begin{enumerate}
  \item 若 $(a,b,c)\neq0$,则
  \[
  \boxed{
  A^+=\dfrac{1}{a^2+b^2+c^2}
  \begin{pmatrix}a\\b\\c\end{pmatrix}
  }.
  \]
  \item 若 $(a,b,c)\neq0$,则
  \[
  \boxed{
  A^+=\dfrac{1}{a^2+b^2+c^2}(a,b,c)
  }.
  \]
  \item 若 $ab\neq0$,则
  \[
  \boxed{A^+=\diag(a^{-1},b^{-1},0)}.
  \]
  \item 若 $(a,b)\neq(0,0)$,写
  \[
  A=uv^{\mathsf T},
  \qquad
  u=\begin{pmatrix}1\\2\end{pmatrix},
  \quad
  v=\begin{pmatrix}a\\b\end{pmatrix}.
  \]
  则
  \[
  \boxed{
  A^+
  =\dfrac{1}{5(a^2+b^2)}
  \begin{pmatrix}
a&2a\\b&2b
\end{pmatrix}
  }.
  \]
  \item 若前两行线性无关,令
  \[
  R=\begin{pmatrix}a&b&c\\e&f&g\end{pmatrix}.
  \]
  则
  \[
  \boxed{
  A^+=R^{\mathsf T}(RR^{\mathsf T})^{-1}(I_2,O)
  }.
  \]
  \item 若前两列线性无关,令
  \[
  C=\begin{pmatrix}a&e\\b&f\\c&g\end{pmatrix}.
  \]
  则
  \[
  \boxed{
  A^+=
  \begin{pmatrix}I_2\\O\end{pmatrix}
  (C^{\mathsf T}C)^{-1}C^{\mathsf T}
  }.
  \]
  \item 若 $(a,b,c)\neq0$,写
  \[
  u=\begin{pmatrix}1\\2\\3\end{pmatrix},
  \qquad
  v=\begin{pmatrix}a\\b\\c\end{pmatrix}.
  \]
  则
  \[
  \boxed{
  A^+
  =\dfrac{1}{14(a^2+b^2+c^2)}
  \begin{pmatrix}
a&2a&3a\\
b&2b&3b\\
c&2c&3c
\end{pmatrix}
  }.
  \]
  \item $J_n(0)$ 是超对角线为 $1$ 的部分置换矩阵,因此
  \[
  \boxed{A^+=J_n(0)^{\mathsf T}}.
  \]
  \item
  \[
  \boxed{E_{1n}^+=E_{n1}}.
  \]
\end{enumerate}
\end{solution}

%例题3.5.5
\begin{example}
证明:
\begin{enumerate}
  \item 对任意矩阵 $A$,矩阵方程
  \[
  AXA=A
  \]
  都有解;
  \item 如果矩阵方程 $AY=C$ 与 $ZB=C$ 都有解,那么方程
  \[
  AXB=C
  \]
  有解.
\end{enumerate}
\end{example}
\exampleblank

\begin{solution}
\leavevmode
\begin{enumerate}
  \item 任意矩阵都存在减号广义逆 $A^-$,取
  \[
  X=A^-
  \]
  即有 $AXA=A$.

  \item 取 $A^-,B^-$ 为任意减号广义逆.由 $AY=C$,
  \[
  AA^-C=AA^-AY=AY=C.
  \]
  由 $ZB=C$,
  \[
  CB^-B=ZBB^-B=ZB=C.
  \]
  令
  \[
  X=A^-CB^-,
  \]
  则
  \[
  AXB=AA^-CB^-B=C.
  \]
\end{enumerate}
\end{solution}

%例题3.5.6
\begin{example}
设 $A$ 为实数域上的 $m\times n$ 矩阵.证明:存在实数域上的 $n\times m$ 矩阵 $B$,使
\[
ABA=A,
\qquad
BAB=B.
\]
并分析:
\begin{enumerate}
  \item 这样的 $B$ 是否唯一?
  \item 若增加要求 $AB$ 与 $BA$ 均为对称矩阵,这样的 $B$ 是否仍然存在?若存在是否唯一?
\end{enumerate}
\end{example}
\exampleblank

\begin{solution}
设 $\rank(A)=r$.取可逆矩阵 $P,Q$ 使
\[
PAQ=E_r=
\begin{pmatrix}I_r&O\\O&O\end{pmatrix}.
\]
于是
\[
A=P^{-1}E_rQ^{-1}.
\]
令
\[
B=QE_r^{\mathsf T}P.
\]
因为
\[
E_rE_r^{\mathsf T}E_r=E_r,
\qquad
E_r^{\mathsf T}E_rE_r^{\mathsf T}=E_r^{\mathsf T},
\]
故
\[
ABA=A,
\qquad
BAB=B.
\]
所以这样的 $B$ 总存在.

一般情况下并不唯一.例如
\[
A=\begin{pmatrix}1&0\\0&0\end{pmatrix}
\]
时,
\[
B=
\begin{pmatrix}
1&s\\t&st
\end{pmatrix},
\qquad s,t\in\mathbb{R},
\]
均满足 $ABA=A,BAB=B$.

若再要求 $AB$ 与 $BA$ 都为对称矩阵,则四个 Penrose 条件全部成立,因此 $B$ 必须是 $A$ 的 Moore--Penrose 广义逆.它总存在且唯一.所以
\[
\boxed{B=A^+}.
\]
\end{solution}
